\documentclass[UTF8,a4paper,11pt]{ctexart}
\usepackage{geometry}
\usepackage{enumitem}
\usepackage{fancyvrb}
\usepackage{booktabs}
\usepackage{hyperref}
\geometry{left=2.2cm,right=2.2cm,top=2.2cm,bottom=2.2cm}
\setlist{nosep}
\hypersetup{colorlinks=true,linkcolor=black,urlcolor=blue}

\title{26-Arch Lab+ 实验报告}
\author{潘孝圆 \quad 24300240128}
\date{2026-06-19}

\begin{document}
\maketitle

\section{完成内容}

本次 Lab+ 在已有 Lab1-Lab6 实现基础上继续补充 bonus。主要完成内容如下：

\begin{itemize}
  \item 保留并验证 Lab1 extra 的乘除法扩展支持。
  \item 保留 Nexys4 上板适配，继续使用非 Basys3 开发板配置。
  \item 保留 Lab5 的 Sv39 MMU、2 MiB/1 GiB hugepage 支持、特权级和 Lab6 异常中断实现。
  \item 新增 Lab+ atomic extension：实现 AMO W 系列指令以及 \texttt{LR.W/SC.W}，并接入 difftest atomic event。
  \item 新增前端性能优化：顺序取指提前发起请求，并为 \texttt{CBusArbiter} 增加 idle fast path，减少固定仲裁空泡。
  \item 新增 PMP/privfull 支持：支持 \texttt{pmpcfg0} 中 8 个 PMP entry 的 \texttt{OFF/TOR/NA4/NAPOT} 匹配，产生 instruction/load/store access fault，并通过 \texttt{lab+/4} privileged sys-test。
  \item 新增 \texttt{EBREAK} 断点异常支持：SYSTEM/funct12=\texttt{0x001} 触发同步异常 cause 3。
  \item 新增 \texttt{FENCE/FENCE.I} 合法 no-op 支持，提升编译器生成程序的兼容性。
  \item 新增 xv6 主线部分进展：补充真实 S-mode、\texttt{SRET}、异常/中断委托和 S 态 trap CSR 写入路径。
  \item 新增 MMU page fault 与 PTE 权限检查：识别 Sv39 非 canonical 地址、无效 PTE、叶子页权限不满足、巨页 PPN 未对齐，并产生 instruction/load/store page fault。
  \item 新增 \texttt{SUM/MXR} 支持：\texttt{MXR} 允许 load 读取 execute-only 页，\texttt{SUM} 允许 S 态数据访问 U 页，同时保持 S 态不能从 U 页取指。
  \item 新增 PTE A/D 位硬件更新：叶子 PTE 的 \texttt{A=0} 或写访问 \texttt{D=0} 时，MMU 先写回更新后的 PTE，再继续最终访存。
  \item 新增 \texttt{test-labplus-2/3/4} 三个 Makefile 测试入口，并补入官方 Lab+ ready-to-run 测试文件。
\end{itemize}

\texttt{lab+/4} 全量 \texttt{TEST=all} 已完成 benchmark 和 sys-test，最终输出 \texttt{Privileged test finished. Exit with code = 0}。当前官方 \texttt{all-test-privfull.bin} 中未包含真实 \texttt{ebreak} 指令，\texttt{breakpoint [X]} 来自测试程序自身的占位输出；补充 \texttt{EBREAK} 后该输出仍不会变化，不影响最终 privileged 测试收尾。

\section{Atomic Extension 设计}

\subsection{指令解码}

新增 opcode \texttt{0101111} 的 A 扩展解码，当前支持 \texttt{funct3=010} 的 32 位原子指令：\texttt{AMOADD.W}、\texttt{AMOSWAP.W}、\texttt{AMOXOR.W}、\texttt{AMOAND.W}、\texttt{AMOOR.W}、\texttt{AMOMIN.W}、\texttt{AMOMAX.W}、\texttt{AMOMINU.W}、\texttt{AMOMAXU.W}、\texttt{LR.W}、\texttt{SC.W}。

\texttt{LR.W} 要求 \texttt{rs2=0}。所有 AMO W 指令按照 store/AMO 地址对齐规则检查地址低两位，不对齐时产生 store/AMO address misaligned 异常。

\subsection{两阶段访存}

当前 CPU 是顺序单发结构，同一时刻只有一条指令处于访存阶段。普通 AMO 指令被实现为两阶段状态机：

\begin{Verbatim}[fontsize=\small]
read old word from memory
compute new word
write new word back to same address
commit rd = sign_extend(old word)
\end{Verbatim}

第一阶段读响应到达后，指令仍保持 \texttt{mem\_pending=1}，并将同一个请求槽切换为写请求。这样不会让取指或下一条指令插入 AMO 的读写之间，满足单核测试中的原子性要求。

\texttt{LR.W} 只执行读请求，并记录 \texttt{\{addr[63:2], 2'b00\}} 作为 reservation 地址。\texttt{SC.W} 在 reservation 命中时发出写请求并返回 0；未命中时不访问内存，直接返回 1，并清除 reservation。

\subsection{Difftest 事件}

AMO 不能直接作为普通 store event 上报。修复方式如下：

\begin{itemize}
  \item 普通 \texttt{DifftestStoreEvent} 屏蔽 \texttt{dt\_is\_atomic}。
  \item AMO 访存提交时额外上报 \texttt{DifftestAtomicEvent}。
  \item AMO load event 的 \texttt{fuType} 设置为 \texttt{0x0f}，使 difftest 按 atomic 路径处理。
  \item \texttt{atomicData} 使用未移位的 \texttt{rs2}，\texttt{atomicOut} 使用旧值，\texttt{atomicMask} 根据地址低三位生成。
\end{itemize}

\section{Privfull 与 PMP 支持}

\subsection{\texttt{jalr} 不对齐异常}

\texttt{lab+/4} 的 \texttt{instr\_misalign} 使用 \texttt{jalr} 跳转到 \texttt{target + 1}。原实现使用清除 bit0 后的 \texttt{jalr} 目标判断是否不对齐，因此该项不会触发异常。为匹配课程测试，本次改为正常跳转目标仍使用 \texttt{\{raw\_target[63:1], 1'b0\}}，但 instruction address misaligned 判断和 \texttt{mtval} 使用未清 bit0 的 \texttt{raw\_target}。该修改后 Lab6 的 \texttt{instr\_misalign} 回归仍通过。

\subsection{PMP 匹配与权限}

实现了 PMP 多 entry 匹配，用于覆盖 Lab+ privfull 测试并补齐更通用的 PMP 行为：

\begin{itemize}
  \item 支持 \texttt{pmpaddr0} 到 \texttt{pmpaddr7}，\texttt{pmpcfg0} 中每个 8-bit 配置项对应一个 PMP entry。
  \item 支持每个 entry 的 \texttt{OFF}、\texttt{TOR}、\texttt{NA4}、\texttt{NAPOT} 地址匹配。
  \item \texttt{TOR} entry 使用前一个 \texttt{pmpaddr} 作为下界，entry0 的下界为 0。
  \item 低编号 entry 优先，命中后使用该 entry 的 \texttt{R/W/X/L} 权限判断。
  \item 所有 PMP entry 都为 \texttt{OFF} 时不影响原有 Lab5/Lab6。
  \item PMP 激活后，U/S 态未命中 PMP 或权限不足会触发 access fault。
  \item M 态在命中 entry 未 lock 时绕过权限检查，命中 lock entry 时按权限检查。
\end{itemize}

测试程序设置的 PMP 范围为 \texttt{0x80006000} 到 \texttt{0x80007000}，即 4096B 用户区。PMP 激活后，用户态访问该范围外的数据或跳转到该范围外取指都会触发异常。

\subsection{异常接入}

取指侧新增 \texttt{if\_id\_access\_fault}。当 PMP 拒绝取指时，CPU 不再真的向内存发起请求，而是向 decode 注入一条携带 fault 信息的占位指令，随后走统一 trap 路径，产生 cause 1。

数据侧在发起 \texttt{dreq} 之前检查 PMP：load/LR/AMO read 权限不足产生 cause 5；store/SC/AMO write 权限不足产生 cause 7；地址不对齐优先级仍高于 access fault。

另外新增 \texttt{fetch\_priv} 解决 trap/mret 重定向同周期的权限判断问题：trap 重定向到 \texttt{mtvec} 时按 M 态检查取指，\texttt{mret} 重定向时按返回后的权限检查取指。

\subsection{\texttt{EBREAK} 断点异常}

Lab+ privfull 输出中存在 \texttt{Test breakpoint [X]}。反汇编和二进制字节检查显示当前官方 \texttt{all-test-privfull.bin} 没有 \texttt{ebreak} 编码 \texttt{0x00100073}，该项不会真正触发断点异常。为补齐异常类型，本次仍在 CPU 中新增 \texttt{EBREAK}：

\begin{itemize}
  \item 在 SYSTEM 指令 \texttt{funct3=000} 且 \texttt{instr[31:20]=12'h001} 时识别为 \texttt{EBREAK}。
  \item \texttt{EBREAK} 作为合法指令进入同步异常路径。
  \item trap cause 设置为 3，\texttt{mtval} 保持 0。
  \item 已重跑 \texttt{make test-labplus-3}、Lab+ privfull sys-test 和 Lab6 sys-test，均保持通过。
\end{itemize}

\subsection{\texttt{FENCE/FENCE.I} 兼容}

当前 CPU 没有 cache，也没有乱序访存或 speculative 执行，因此普通 \texttt{FENCE} 和 \texttt{FENCE.I} 不需要真实刷新硬件状态。本次将 opcode \texttt{0001111} 中 \texttt{funct3=000/001} 的指令作为合法 no-op 提交：

\begin{itemize}
  \item \texttt{FENCE} 用于兼容内存顺序屏障。
  \item \texttt{FENCE.I} 用于兼容指令流同步屏障。
  \item 两者不写寄存器、不发访存请求、不触发流水线重定向。
  \item 已随 atomic、Lab+ privfull 和 Lab6 回归一起验证。
\end{itemize}

\subsection{xv6 主线部分进展：S-mode 与委托}

官方 Lab+ 主 Track 是尝试运行 xv6。当前仓库没有 xv6 镜像或源码，因此本次先补 xv6 boot 的 CPU 前置能力，而不是直接实现磁盘 MMIO。新增内容如下：

\begin{itemize}
  \item 新增 \texttt{PRIV\_S=2'b01}，\texttt{mret} 可以返回到 S 态，difftest 的 privilege mode 也能看到 S 编码。
  \item 新增 \texttt{SRET} 解码，S/M 态可执行，返回到 \texttt{sstatus.spp} 指定的 U/S 态，并按规范更新 \texttt{sstatus.sie/spie/spp}。
  \item \texttt{ECALL} 在 S 态触发 cause 9，U 态仍为 cause 8，M 态仍为 cause 11。
  \item 开放 \texttt{medeleg} 常见可委托异常位，开放 \texttt{mideleg} 的 SSIP/STIP/SEIP 位。
  \item trap 被委托到 S 态时，跳转 \texttt{stvec}，写入 \texttt{sepc/scause/stval}，并更新 \texttt{sstatus.sie/spie/spp}；未委托时保持原 M 态 trap 路径。
  \item 增加 S 级中断 evaluate 框架：当 \texttt{mip/sip}、\texttt{mie/sie}、\texttt{mideleg} 同时打开时，U/S 态可进入 S trap。
\end{itemize}

这一部分还不能直接跑完整 xv6，因为仍缺少 virtio/磁盘 MMIO 或替代块设备模型。但它把 xv6 所需的 Supervisor trap 基础路径补上了，并通过现有 Lab5、Lab6 与 Lab+ privfull 回归确认没有破坏原 U/M 行为。

\subsection{MMU page fault 与 PTE 权限检查}

原 MMU 在页表项无效时会把最终物理地址置 0 并继续发起访存，这对 Lab5 的正常路径足够，但不适合继续做 xv6/异常测试。本次在 CBus 响应中增加 \texttt{page\_fault} 标记，并把该标记一路传回取指和数据访存路径：

\begin{itemize}
  \item CBus 请求增加 \texttt{is\_instr}，MMU 能区分取指、load 读和 store/AMO 写。
  \item CBus/IBus/DBus 响应增加 \texttt{page\_fault}，MMU 发现翻译失败时返回一拍 fault 响应，不再访问物理地址 0。
  \item 取指 page fault 进入 decode 统一异常路径，产生 cause 12，\texttt{mtval/stval} 写入 faulting virtual address。
  \item load/LR/AMO read page fault 在数据响应阶段产生 cause 13。
  \item store/SC/AMO write page fault 在数据响应阶段产生 cause 15，并阻止该访存提交到 difftest。
  \item MMU 检查 Sv39 canonical 地址、\texttt{V=0}、\texttt{W=1 \&\& R=0}、第三级仍非叶子、叶子页 \texttt{R/W/X/U} 权限和 1 GiB/2 MiB 巨页 PPN 对齐。
\end{itemize}

数据侧 page fault 是响应期异常，不在 decode 阶段就能确定。因此本次额外补了 WB trap 重定向：当 \texttt{dresp.page\_fault=1} 时清除等待中的访存、写入 trap CSR，并清空 IF/ID，防止后续指令越过 faulting load/store 提交。

\subsection{PTE A/D 位硬件更新}

RISC-V Sv39 的叶子 PTE 中，\texttt{A} 表示 accessed，\texttt{D} 表示 dirty。为了让不预先置 A/D 的页表也可以运行，本次在 MMU 页表 walker 中加入硬件更新路径：

\begin{itemize}
  \item 叶子 PTE 权限检查通过后，如果 \texttt{A=0}，MMU 对该 PTE 地址写回 \texttt{A=1}。
  \item 对 store/SC/AMO write，如果 \texttt{D=0}，同一次写回同时置 \texttt{D=1}。
  \item instruction fetch 和 load 只需要置 \texttt{A}，不置 \texttt{D}。
  \item PTE 写回完成后，MMU 再发起原本的最终物理地址访问。
  \item 非法 PTE、权限不足、superpage PPN 未对齐等 fault 不会更新 A/D 位。
\end{itemize}

实现上新增 \texttt{S\_AD\_UPDATE} 状态，复用已有 CBus 对页表项地址发起 64 位写事务；外部 MMU 接口保持不变。

\subsection{\texttt{SUM/MXR} 权限补充}

在 \texttt{MMU} 中继续补充了 \texttt{mstatus.sum} 和 \texttt{mstatus.mxr} 对页表权限的影响：

\begin{itemize}
  \item \texttt{mstatus.mxr=1} 时，load 可以读取 \texttt{X=1/R=0} 的 execute-only 页。
  \item \texttt{mstatus.mxr=0} 时，load 仍要求 \texttt{R=1}。
  \item U 态访问页表项时仍要求 \texttt{PTE\_U=1}。
  \item S 态访问 supervisor 页时要求 \texttt{PTE\_U=0}。
  \item S 态数据访问用户页时，只有 \texttt{mstatus.sum=1} 才允许。
  \item S 态取指仍然禁止从用户页执行，即使 \texttt{SUM=1} 也会触发 instruction page fault。
\end{itemize}

实现上，\texttt{core} 将当前 \texttt{mstatus} 输出到 MMU，MMU 在叶子 PTE 权限判断时组合使用 \texttt{priv\_mode}、\texttt{is\_instr}、\texttt{is\_write}、\texttt{SUM/MXR} 和 \texttt{R/W/X/U} 位。该改动不改变 CSR 写入路径，\texttt{MSTATUS\_MASK/SSTATUS\_MASK} 之前已经开放了 \texttt{SUM/MXR} 位，因此软件可以通过 \texttt{mstatus/sstatus} 正常设置。

\section{前端性能优化}

\subsection{顺序取指提前}

优化前取指状态机只在以下条件满足时发起新请求：

\begin{Verbatim}[fontsize=\small]
!if_pending && !if_id_valid && !mem_pending
\end{Verbatim}

优化后新增 \texttt{if\_can\_request}：

\begin{Verbatim}[fontsize=\small]
!if_pending && !mem_pending
&& (!if_id_valid || (id_consume && !id_redirect))
\end{Verbatim}

当当前 \texttt{if\_id} 指令在本周期被消费，并且不是 trap、跳转、taken branch、CSR 或 MRET 时，CPU 同周期发起下一条顺序取指请求。对于重定向类指令，仍然等待目标写入 \texttt{pc} 后再取指。

\subsection{\texttt{CBusArbiter} fast path}

原 \texttt{CBusArbiter} 在 idle 状态下不会直接把选中的请求透传到输出总线，因此每个请求都会固定多等一拍。本次改为 idle 且存在请求时，\texttt{oreq} 直接等于当前选中的 \texttt{selected\_req}；如果下游同周期返回 \texttt{last}，仲裁器不进入 busy；如果请求未完成，才锁存 \texttt{index} 并进入 busy。该优化对取指和数据访存都生效。

\subsection{周期对比}

\begin{center}
\begin{tabular}{lrrrrr}
\toprule
测试 & 原始周期 & 顺序提前 & 加 fast path & 总减少 & 最终 IPC \\
\midrule
\texttt{atomicity.bin} & 372 & 317 & 249 & 33.1\% & 0.221 \\
\texttt{lab1-extra-test.bin} & 185976 & 153201 & 120425 & 35.2\% & 0.272 \\
\texttt{lab4-test.bin} & 208529 & 177990 & 139050 & 33.3\% & 0.236 \\
\bottomrule
\end{tabular}
\end{center}

MicroBench qsort 采样从 7 ms 降到 6 ms。\texttt{lab+/4 TEST=all} 中的 benchmark 结果如下：

\begin{Verbatim}[fontsize=\small]
CoreMark: 738 ms, 13 Iterations/Sec
Dhrystone: 1046 ms, 16 Marks
STREAM Copy/Scale/Add/Triad: 14.5 / 0.9 / 1.8 / 0.8 MB/s
\end{Verbatim}

\section{代码修改}

\begin{itemize}
  \item \texttt{vsrc/src/core.sv}：新增 AMO 解码、AMO 结果计算、reservation 记录；扩展访存状态机；新增 difftest atomic event；新增 \texttt{if\_can\_request}；新增 8-entry PMP 匹配、取指 access fault 注入、load/store access fault；新增 \texttt{fetch\_priv}；新增 \texttt{EBREAK} 解码和 breakpoint exception cause 3；新增 \texttt{FENCE/FENCE.I} 合法 no-op 解码；新增 S-mode、\texttt{SRET}、\texttt{medeleg/mideleg} 委托、S 态 trap CSR 写入；新增 instruction/load/store page fault 接入和数据访存 fault 后的 WB trap 清流水；输出当前 \texttt{mstatus} 给 MMU 使用。
  \item \texttt{vsrc/include/common.sv}：扩展 CBus/IBus/DBus 响应，增加 \texttt{page\_fault}；扩展 CBus 请求，增加 \texttt{is\_instr} 标记。
  \item \texttt{vsrc/util/MMU.sv}：新增 fault 状态和 PTE 权限检查，不再将无效 PTE 转成物理地址 0；支持 Sv39 canonical 地址检查和巨页 PPN 对齐检查；支持 \texttt{SUM/MXR} 对 S/U 态数据访问和 execute-only 页读取的影响；新增 \texttt{S\_AD\_UPDATE} 状态，支持硬件置 PTE \texttt{A/D} 位。
  \item \texttt{vsrc/util/IBusToCBus.sv}、\texttt{vsrc/util/DBusToCBus.sv}：传递 \texttt{is\_instr} 与 \texttt{page\_fault}。
  \item \texttt{vsrc/mycpu\_top.sv}、\texttt{vsrc/util/CBusToAXI.sv}、\texttt{vivado/src/with\_delay/soc\_top.sv}：对真实内存响应源固定 \texttt{page\_fault=0}。
  \item \texttt{vsrc/include/csr.sv}：修正 \texttt{SSTATUS\_MASK}，开放 SIE/SPIE/SPP 等 S 态状态位；设置 \texttt{MEDELEG\_MASK/MIDELEG\_MASK}，支持常见异常和 S 级中断委托；新增 \texttt{pmpaddr1} 到 \texttt{pmpaddr7} CSR 地址常量。
  \item \texttt{vsrc/util/CBusArbiter.sv}：idle 状态下直接透传当前请求，同周期完成的请求不再进入 busy。
  \item \texttt{Makefile}：新增 \texttt{test-labplus-2}、\texttt{test-labplus-3}、\texttt{test-labplus-4}。
  \item \texttt{ready-to-run/lab+/}：补充官方 Lab+ 测试二进制和汇编反汇编文件。
\end{itemize}

\section{测试结果}

\subsection{Lab+ atomic extension}

\begin{Verbatim}[fontsize=\small]
make test-labplus-3
The image is ./ready-to-run/lab+/3/atomicity.bin
Core 0: HIT GOOD TRAP at pc = 0x800000dc
total guest instructions = 55
instrCnt = 55, cycleCnt = 249
\end{Verbatim}

\subsection{Lab+ privfull/PMP}

\begin{Verbatim}[fontsize=\small]
Test ecall_u [OK]
Test instr_misalign [OK]
Test instr_access_fault [OK]
Test illegal_instr [OK]
Test breakpoint [X]
Test load_misalign [OK]
Test load_fault [OK]
Test store_misalign [OK]
Test store_fault [OK]
Test timer_intr [OK]
Test software_intr [OK]
Test pmp_nr [OK]
Test pmp_nw [OK]
Test pmp_nx [OK]
Test mem_detect [OK]
    4096B user memory detected. (Interrupts: 27899)
Test m_trap [OK]
Privileged test finished.
Exit with code = 0
\end{Verbatim}

\subsection{回归测试}

\begin{Verbatim}[fontsize=\small]
Lab1 extra:
Core 0: HIT GOOD TRAP at pc = 0x8002001c
instrCnt = 32775, cycleCnt = 120425

Lab4:
Core 0: HIT GOOD TRAP at pc = 0x8001fff8
instrCnt = 32766, cycleCnt = 139050

Lab5:
xv6 kernel is booting
Return from init! Test passed

Lab6:
Test ecall_u [OK]
Test instr_misalign [OK]
Test load_misalign [OK]
Test store_misalign [OK]
Test timer_intr [OK]
Test software_intr [OK]
Test m_trap [OK]
Privileged test finished.
Exit with code = 0
\end{Verbatim}

\subsection{S-mode/SRET 回归}

本次新增 S-mode、委托路径、MMU page fault、PTE 权限检查、\texttt{SUM/MXR}、PTE A/D 硬件更新与 8-entry PMP 后重新运行：

\begin{Verbatim}[fontsize=\small]
make test-labplus-3
Core 0: HIT GOOD TRAP at pc = 0x800000dc
instrCnt = 55, cycleCnt = 249

Lab5 kernel:
Return from init! Test passed

TEST=sys ./build/emu --no-diff -i ./ready-to-run/lab6/lab6-test.bin
Privileged test finished.
Exit with code = 0

TEST=sys ./build/emu --no-diff -i ./ready-to-run/lab+/4/all-test-privfull.bin
Privileged test finished.
Exit with code = 0
\end{Verbatim}

\section{后续可做项}

本次已经完成 xv6 主线的前两步：S-mode/trap delegation，以及 MMU page fault/PTE 基础权限检查。后续如果继续推进 xv6，需要补充更有针对性的 page fault 单元测试、virtio 或简化磁盘 MMIO、S 态 timer/external interrupt 与 CLINT/PLIC 的转换路径。

性能方向还可以继续实现 cache、分支预测或多级流水线，用于进一步提升 \texttt{test-labplus-2}。这些改动会触及 trap、CSR、MMU、取指和访存多个共享路径，风险明显高于本次 8-entry PMP、atomic 和前端空泡优化，因此本次没有继续扩大范围。

\section{总结}

本次 Lab+ 新增完成了 atomic extension、8-entry PMP/privfull 支持、\texttt{EBREAK} 断点异常、\texttt{FENCE/FENCE.I} 兼容，并加入两项轻量性能优化。继续推进 xv6 主 Track 时，已完成 S-mode、\texttt{SRET}、trap delegation、MMU page fault/PTE 基础权限检查、\texttt{SUM/MXR} 权限补充，以及 PTE A/D 位硬件更新。AMO 实现利用现有单发访存结构，将 AMO 指令拆成不可被其他指令插入的读-改-写序列，并为 \texttt{LR.W/SC.W} 添加 reservation 状态。性能优化将 Lab1 extra 周期数从 185976 降到 120425，将 Lab4 周期数从 208529 降到 139050。最终 atomic、Lab+ privileged sys-test、Lab1 extra、Lab4、Lab5、Lab6 均完成回归。

\section{AI 使用说明}

我作为本项目的主导人，负责实验方案制定、代码审阅、测试验证和最终提交确认。AI（Codex）用于辅助阅读 Lab+ 要求、分析官方 \texttt{atomicity.S} 和 \texttt{all-test-privfull.S}、实现 AMO/PMP/性能优化、定位 difftest atomic event 与 trap 重定向时序问题，并整理实验报告。最终代码和报告由我本人审阅确认。

\end{document}
